% !Mode:: "TeX:UTF-8"
% !TEX root = ../../main.tex
% ============================================================
% TOPIC 16
% ============================================================
\mashgulot{16}{Definition of a Euclidean space. The norm of an element of a Euclidean space}

\begin{tarif}
Let $E$ be a real linear space, and suppose that to every pair of elements
$x,y\in E$ a real number $(x,y)$ is assigned so that the following axioms hold.
Then $(x,y)$ is called an \textbf{inner (scalar) product}, and $E$ is called a
\textbf{Euclidean space}:
\begin{enumerate}[label=\arabic*)]
\item $(x,y)=(y,x)$ (symmetry);
\item $(x+y,z)=(x,z)+(y,z)$; \quad $(\lambda x,y)=\lambda(x,y)$ (linearity);
\item $(x,x)\geq0$, and $(x,x)=0\iff x=\theta$ (positivity).
\end{enumerate}
\end{tarif}

\textbf{Examples.} 1) The standard inner product on $\R^n$ is
$(x,y)=x_1y_1+x_2y_2+\dots+x_ny_n$.
2) On $\R^2$ one can take the weighted product $(x,y)=2x_1y_1+3x_2y_2$.
3) On $C[a;b]$ the formula $(f,g)=\displaystyle\int_a^b f(t)g(t)\,dt$ defines an inner product.

\subsection{Norm, distance and angle}

\begin{tarif}
The \textbf{norm (length)} of an element $x$ is $\|x\|=\sqrt{(x,x)}$; the
\textbf{distance} between two elements is $\rho(x,y)=\|x-y\|$; the angle
$\varphi$ between them is defined by
\[\kf{\cos\varphi=\frac{(x,y)}{\|x\|\,\|y\|}\qquad(x,y\neq\theta).}\]
If $(x,y)=0$, the elements $x$ and $y$ are said to be \textbf{orthogonal}, and we write $x\perp y$.
\end{tarif}

\begin{teorema}
In every Euclidean space:
\begin{enumerate}[label=\arabic*)]
\item $|(x,y)|\leq\|x\|\,\|y\|$ (the \textbf{Cauchy–Schwarz} or \textbf{Cauchy–Bunyakovsky–Schwarz inequality});
\item $\|x+y\|\leq\|x\|+\|y\|$ (the \textbf{triangle inequality});
\item if $x\perp y$, then $\|x+y\|^2=\|x\|^2+\|y\|^2$ (the \textbf{Pythagorean theorem}).
\end{enumerate}
\end{teorema}

\begin{misol}
In $\R^4$ let $x=(1;2;-1;3)$ and $y=(2;0;1;1)$. Find $(x,y)$, $\|x\|$, $\|y\|$,
$\cos\varphi$ and $\rho(x,y)$.

\textbf{Solution.} We have $(x,y)=2+0-1+3=4$, \ $\|x\|=\sqrt{1+4+1+9}=\sqrt{15}$ and
$\|y\|=\sqrt{4+0+1+1}=\sqrt6$; hence
\[\kf{\cos\varphi=\frac{4}{\sqrt{15}\sqrt6}=\frac{4}{3\sqrt{10}},\qquad
\rho(x,y)=\|(-1;2;-2;2)\|=\sqrt{13}.}\]
The Cauchy–Schwarz inequality is satisfied, as it must be: $4\leq\sqrt{90}\approx9.49$.
\end{misol}

\begin{misol}
Are the vectors $x=(1;2)$ and $y=(3;-1)$ orthogonal with respect to the inner
product $(x,y)=2x_1y_1+3x_2y_2$ on $\R^2$?

\textbf{Solution.} Here $(x,y)=2\cdot1\cdot3+3\cdot2\cdot(-1)=0$, so the vectors are
orthogonal. With the standard inner product, however, $(x,y)=3-2=1\neq0$.
Thus orthogonality depends on the choice of inner product.
\end{misol}

\begin{misol}
In $C[0;1]$ take $f(t)=t$ and $g(t)=1$. Then
$(f,g)=\int_0^1t\,dt=\frac12$, \ $\|f\|=\sqrt{\int_0^1t^2dt}=\frac{1}{\sqrt3}$, \
$\|g\|=1$, and $\cos\varphi=\frac{1/2}{1/\sqrt3}=\frac{\sqrt3}{2}$, so that $\varphi=30^\circ$.
\end{misol}

\begin{tarif}
To \textbf{normalize} an element $x\neq\theta$ means to pass to the unit element
$x^0=\dfrac{x}{\|x\|}$.
\end{tarif}

\subsection{Applications of the norm in information security}

\textbf{Hamming distance.} In the space of binary words, $d(u,v)$ is the number of
positions in which $u$ and $v$ differ; it equals the number of ones in $u\oplus v$.
If a code has minimum distance $d_{\min}$, it detects up to $d_{\min}-1$ errors
and corrects up to $\left\lfloor\frac{d_{\min}-1}{2}\right\rfloor$ errors.

\begin{misol}
For $u=1011010$ and $v=1001110$ we get $u\oplus v=0010100$, hence $d(u,v)=2$.
\end{misol}

\textbf{Anomaly detection.} A user’s activity is described by a vector $\vec x$;
if its Euclidean distance $\|\vec x-\vec m\|$ from the normal profile $\vec m$
exceeds a given threshold, the behaviour is flagged as suspicious. In biometric
authentication (face, fingerprint) templates are also stored as vectors, and
they are compared by means of a distance or the cosine similarity.

\textbf{Lattice cryptography.} The shortest vector problem (SVP) asks for a
nonzero lattice vector of smallest norm. In high dimensions it is extremely hard
computationally, which is why it serves as one of the basic hard problems of
post-quantum cryptography.

\subsection{Computing with Python, Maple and Mathcad}

The quantities of the first example above can be obtained in a few lines. In
Python we use SymPy, which keeps the square roots exact:
\begin{Verbatim}[fontsize=\footnotesize,frame=single,label={\small Python}]
from sympy import Matrix, radsimp
x = Matrix([1, 2, -1, 3])
y = Matrix([2, 0, 1, 1])
print(x.dot(y), x.norm(), y.norm())
print(radsimp(x.dot(y)/(x.norm()*y.norm())))
print((x - y).norm())
\end{Verbatim}
\begin{Verbatim}[fontsize=\footnotesize,frame=single,label={\small Output}]
4 sqrt(15) sqrt(6)
2*sqrt(10)/15
sqrt(13)
\end{Verbatim}
Note that $\frac{2\sqrt{10}}{15}=\frac{4}{3\sqrt{10}}$, in agreement with the hand computation.
In Maple the same results are given by the \texttt{LinearAlgebra} package
(\texttt{VectorAngle} returns the angle itself):
\begin{Verbatim}[fontsize=\footnotesize,frame=single,label={\small Maple}]
with(LinearAlgebra):
x := <1, 2, -1, 3>:  y := <2, 0, 1, 1>:
DotProduct(x, y);              # 4
Norm(x, 2);  Norm(y, 2);       # sqrt(15), sqrt(6)
DotProduct(x, y)/(Norm(x, 2)*Norm(y, 2));
Norm(x - y, 2);                # sqrt(13)
VectorAngle(x, y);
\end{Verbatim}
In Mathcad the vectors are entered as columns; the dot product is the
multiplication operator applied to two vectors, and $|x|$ is the length of a vector:
\begin{Verbatim}[fontsize=\footnotesize,frame=single,label={\small Mathcad}]
     | 1 |        | 2 |
x := | 2 |   y := | 0 |
     |-1 |        | 1 |
     | 3 |        | 1 |
x.y = 4     |x| = 3.873     |y| = 2.449
  x.y
-------- = 0.422     |x - y| = 3.606
|x|.|y|
\end{Verbatim}

\subsection{Exercises on Euclidean spaces and norms}

\emph{1. Find the inner product and the norms of the vectors.}
\bigskip
\begin{enumerate}[label=\arabic*., leftmargin=*]
\item $x=(0;\,2;\,2;\,-1)$, $y=(1;\,-2;\,2;\,-5)$
\item $x=(1;\,5;\,-1;\,-2)$, $y=(5;\,-2;\,-5;\,-1)$
\item $x=(-1;\,0;\,5;\,-3)$, $y=(4;\,-1;\,-5;\,-2)$
\item $x=(4;\,-1;\,-5;\,-3)$, $y=(4;\,5;\,5;\,-5)$
\item $x=(2;\,2;\,4;\,5)$, $y=(-1;\,-2;\,-1;\,0)$
\item $x=(-1;\,1;\,-4;\,0)$, $y=(2;\,1;\,3;\,5)$
\item $x=(-3;\,4;\,-1;\,4)$, $y=(-5;\,-1;\,-4;\,-5)$
\item $x=(3;\,0;\,-2;\,2)$, $y=(-3;\,-1;\,-1;\,0)$
\item $x=(2;\,2;\,-4;\,-3)$, $y=(2;\,-5;\,2;\,2)$
\item $x=(-5;\,2;\,-4;\,2)$, $y=(4;\,-4;\,2;\,5)$
\item $x=(-5;\,-3;\,-2;\,1)$, $y=(0;\,-5;\,3;\,-5)$
\item $x=(5;\,5;\,1;\,4)$, $y=(0;\,2;\,3;\,5)$
\item $x=(5;\,5;\,-4;\,-2)$, $y=(0;\,-4;\,-4;\,3)$
\item $x=(-4;\,-2;\,-5;\,1)$, $y=(5;\,-5;\,-4;\,5)$
\item $x=(-5;\,1;\,-3;\,5)$, $y=(-1;\,-2;\,-3;\,4)$
\item $x=(-5;\,-2;\,3;\,-2)$, $y=(-4;\,-4;\,5;\,-3)$
\item $x=(-1;\,1;\,1;\,-5)$, $y=(1;\,2;\,2;\,2)$
\item $x=(-1;\,-5;\,-5;\,-5)$, $y=(-2;\,1;\,5;\,-1)$
\item $x=(3;\,1;\,-4;\,-2)$, $y=(-4;\,1;\,5;\,-4)$
\item $x=(-1;\,-4;\,2;\,1)$, $y=(-3;\,1;\,-2;\,-3)$
\item $x=(-3;\,1;\,1;\,3)$, $y=(-2;\,1;\,-3;\,0)$
\item $x=(5;\,2;\,-3;\,1)$, $y=(-2;\,5;\,1;\,-5)$
\item $x=(2;\,-1;\,4;\,-1)$, $y=(-4;\,-4;\,-1;\,2)$
\item $x=(-2;\,3;\,5;\,2)$, $y=(-2;\,4;\,2;\,3)$
\item $x=(-3;\,2;\,3;\,-4)$, $y=(5;\,5;\,-4;\,5)$
\item $x=(0;\,-5;\,2;\,-4)$, $y=(2;\,-5;\,3;\,0)$
\item $x=(3;\,-3;\,-4;\,0)$, $y=(1;\,-1;\,1;\,-4)$
\item $x=(-2;\,-2;\,2;\,-2)$, $y=(-5;\,5;\,-5;\,-2)$
\item $x=(0;\,1;\,4;\,4)$, $y=(-2;\,-5;\,-1;\,5)$
\item $x=(-5;\,1;\,5;\,-5)$, $y=(-3;\,5;\,3;\,-2)$
\end{enumerate}

\emph{2. Find the cosine of the angle between the vectors and the distance between them.}
\bigskip
\begin{enumerate}[label=\arabic*., leftmargin=*]
\item $x=(-2;\,3;\,5;\,3)$, $y=(5;\,5;\,-1;\,3)$
\item $x=(-5;\,4;\,4;\,3)$, $y=(-2;\,2;\,4;\,2)$
\item $x=(0;\,-3;\,2;\,-3)$, $y=(-4;\,-2;\,-5;\,-4)$
\item $x=(0;\,-1;\,5;\,5)$, $y=(-4;\,-3;\,-5;\,-3)$
\item $x=(4;\,-4;\,-2;\,0)$, $y=(-2;\,4;\,1;\,-4)$
\item $x=(-3;\,-4;\,2;\,-4)$, $y=(-5;\,4;\,-3;\,-3)$
\item $x=(5;\,-3;\,-2;\,1)$, $y=(-5;\,0;\,-4;\,5)$
\item $x=(1;\,1;\,-1;\,-4)$, $y=(4;\,1;\,4;\,-2)$
\item $x=(0;\,2;\,0;\,1)$, $y=(-5;\,0;\,-2;\,-4)$
\item $x=(-2;\,1;\,1;\,-5)$, $y=(4;\,-2;\,1;\,5)$
\item $x=(3;\,0;\,-5;\,3)$, $y=(1;\,0;\,-5;\,5)$
\item $x=(-3;\,1;\,-2;\,0)$, $y=(-3;\,1;\,-5;\,1)$
\item $x=(-5;\,1;\,5;\,2)$, $y=(5;\,-5;\,-3;\,-5)$
\item $x=(-3;\,-2;\,5;\,4)$, $y=(-1;\,4;\,5;\,4)$
\item $x=(1;\,-5;\,-5;\,-4)$, $y=(3;\,3;\,-5;\,-3)$
\item $x=(0;\,4;\,4;\,-4)$, $y=(4;\,3;\,-1;\,2)$
\item $x=(-3;\,4;\,1;\,-1)$, $y=(4;\,4;\,1;\,3)$
\item $x=(-4;\,-3;\,4;\,-5)$, $y=(-1;\,-1;\,-1;\,5)$
\item $x=(2;\,0;\,-4;\,5)$, $y=(-1;\,1;\,-1;\,5)$
\item $x=(4;\,2;\,-5;\,-4)$, $y=(-4;\,0;\,1;\,-1)$
\item $x=(2;\,-1;\,-2;\,-1)$, $y=(-5;\,2;\,-5;\,-2)$
\item $x=(-4;\,1;\,-3;\,4)$, $y=(-2;\,0;\,1;\,4)$
\item $x=(-4;\,-1;\,-5;\,0)$, $y=(-3;\,0;\,1;\,4)$
\item $x=(-1;\,0;\,-2;\,-3)$, $y=(5;\,-2;\,3;\,-5)$
\item $x=(3;\,-4;\,-1;\,-5)$, $y=(-5;\,3;\,-4;\,-1)$
\item $x=(-2;\,-4;\,-4;\,2)$, $y=(2;\,0;\,-5;\,-5)$
\item $x=(-4;\,-1;\,5;\,1)$, $y=(-4;\,0;\,1;\,4)$
\item $x=(2;\,1;\,-2;\,4)$, $y=(-2;\,-1;\,-3;\,-2)$
\item $x=(5;\,-3;\,-3;\,1)$, $y=(0;\,-4;\,-3;\,-5)$
\item $x=(5;\,-2;\,-3;\,5)$, $y=(1;\,-4;\,-1;\,-1)$
\end{enumerate}

\emph{3. Applied problems (10 problems).}
\bigskip
\begin{enumerate}[label=\arabic*., leftmargin=*]
\item \textbf{(Orthogonality.)} Are the vectors $x=(1;2;-1;3)$ and $y=(2;1;1;-1)$ of $\mathbb R^4$ orthogonal?
\item \textbf{(Normalization.)} Normalize the vector $x=(1;-2;2;4)$.
\item \textbf{(Cauchy–Schwarz.)} Verify the inequality $|(x,y)|\le\|x\|\,\|y\|$ for $x=(1;2;3)$ and $y=(4;-1;2)$.
\item \textbf{(Triangle inequality.)} Verify the inequality $\|x+y\|\le\|x\|+\|y\|$ for $x=(3;0;4)$ and $y=(0;5;12)$.
\item \textbf{(A parameter.)} For which value of $\lambda$ are the vectors $x=(1;\lambda;2;-1)$ and $y=(3;2;\lambda;1)$ orthogonal?
\item \textbf{(Weighted inner product.)} On $\mathbb R^2$ the inner product $(x,y)=2x_1y_1+3x_2y_2$ is given. Find $(x,y)$, $\|x\|$ and $\|y\|$ for $x=(1;2)$ and $y=(3;-1)$.
\item \textbf{(A function space.)} On $C[0;1]$ let $(f,g)=\int_0^1 f(t)g(t)\,dt$. Find $(f,g)$ and $\|f\|$ for $f=t$ and $g=1$.
\item \textbf{(Hamming distance.)} Find the Hamming distance between the binary words $u=1011010$ and $v=1001110$.
\item \textbf{(Cosine similarity.)} The word-frequency vectors of two documents are $u=(3;0;2;1)$ and $v=(1;1;2;0)$. Find $\cos(u,v)$.
\item \textbf{(Pythagorean theorem.)} Verify that $\|x+y\|^2=\|x\|^2+\|y\|^2$ for $x=(1;1;0;0)$ and $y=(0;0;2;2)$.
\end{enumerate}

